% A espaçonave Dragon 2, cujo desenvolvimento foi iniciado no primeiro semestre de 2013 pela SpaceX \cite{spaceDragonAlien}, é projetada para transportar tanto cargas quanto tripulações humanas ao espaço.
% Essa versão se baseia no modelo original \emph{Dragon 1} e apresenta melhorias em termos de autonomia, segurança e capacidade operacional.
% A \emph{Dragon 2} pode transportar até sete tripulantes ou uma combinação de cargas e tripulantes, podendo ser utilizada em diferentes tipos de missões espaciais, como o transporte para a ISS ou a entrega de novos astronautas na ISS para que outros retornem para a Terra.

% Entre as principais melhorias em relação à \emph{Dragon 1}, está o sistema de lançamento e retorno. A \emph{Dragon 2} é equipada com motores \emph{SuperDraco}, que oferecem capacidade de abortagem de emergência durante o lançamento.
% Esse sistema permite que a espaçonave escape em caso de falhas no foguete, protegendo a tripulação. Além disso, a \emph{Dragon 2} utiliza paraquedas redesenhados para garantir um pouso mais controlado, uma atualização em relação ao modelo anterior, que dependia de procedimentos diferentes para retornar à Terra.
% Outra inovação é o sistema de navegação autônoma \cite{nasaStarshipDocking}, que realiza o acoplamento direto à ISS sem a necessidade do braço robótico \emph{Canadarm2}, que em sua versão anterior (\emph{Dragon 1}) exigia para realizar o acoplamento na ISS.
% Essa funcionalidade reduz melhora a eficiência das missões pois eleva a segurança, reduz a dependência humana e pode ocorrer mesmo que haja restrições como ausência de iluminação ou tripulação em horário de descanso.
% Por fim, a \emph{Dragon 2} apresenta um design interno atualizado, com controles interativos por \emph{touchscreen} e maior conforto para longas viagens.

%%%%%%%%%%%%%%%%%%%%%%%%%%%%%%%%

A espaçonave \emph{Dragon 2}, cujo desenvolvimento foi iniciado no primeiro semestre de 2013 pela SpaceX \cite{spaceDragonAlien}, foi projetada para transportar carga e tripulação.
Seu projeto foi derivado da espaçonave anterior, \emph{Dragon 1}, e incorpora mudanças voltadas a maior grau de autonomia, requisitos de segurança e ampliação de capacidade operacional.
Pode transportar até sete tripulantes ou uma configuração mista de carga e tripulação, em perfis de missão como transporte e rotação de astronautas na ISS.

Em relação à \emph{Dragon 1}, destacam-se alterações nos sistemas de lançamento, contingência e retorno. A \emph{Dragon 2} é equipada com motores \emph{SuperDraco}, que permitem abortagem de emergência durante a ascensão, aumentando a proteção da tripulação em cenários de falha do lançador. No retorno, utiliza paraquedas redesenhados para controlar o pouso.
Outra diferença relevante é a capacidade de navegação e acoplamento autônomos \cite{nasaStarshipDocking}, com acoplamento direto à ISS sem uso do braço robótico \emph{Canadarm2}, ao contrário do perfil típico da \emph{Dragon 1}.
Por fim, o interior foi atualizado com interfaces por \emph{touchscreen} e ajustes ergonômicos para missões de maior duração.
